% ============================================================
% Appendix B: IBM Quantum Hardware Experiments
% Tommaso R. Marena (2026)
%
% Fill in the [BRACKETED] placeholders after running
% notebooks/hardware_ibm_colab.ipynb and
% notebooks/circuit_depth_scaling.ipynb.
% ============================================================

\section*{Appendix B: Hardware Validation}
\label{app:hardware}

\subsection*{B.1 Experimental Setup}

We validate the QOS Boolean phase-oracle circuit on real superconducting
quantum hardware via the IBM Quantum platform~\cite{ibmquantum2024}.
All experiments were conducted on \textbf{[BACKEND\_NAME]} ([BACKEND\_QUBITS]
physical qubits) using Qiskit Runtime v[QISKIT\_RUNTIME\_VERSION].
Circuits were compiled with optimization level~3 and executed with
[SHOTS] measurement shots per circuit.

For each circuit size $n \in \{3, 4, 5\}$ we prepared a QOS Boolean oracle
for a uniformly random $n$-bit truth table with sample budget
$M = [M\_SAMPLES]$, matching the regime studied in Section~6.
The oracle diagonal was decomposed into hardware-native gates via
a Walsh--Hadamard Transform (WHT) decomposition (see Appendix~A).

\subsection*{B.2 Zero-Noise Extrapolation}

To mitigate the effect of gate noise, we applied
Zero-Noise Extrapolation (ZNE) via gate-folding~\cite{temme2017error,li2017efficient}.
Each two-qubit gate $G$ was replaced by $G\,(G^\dagger G)^{(\lambda-1)/2}$ for
noise scale factors $\lambda \in \{[ZNE\_FACTORS]\}$,
and the resulting probability distributions were fit by a degree-$[ZNE\_DEG]$
polynomial and extrapolated to $\lambda = 0$.

\subsection*{B.3 Results}

Table~\ref{tab:hardware} reports the classical fidelity and total variation
distance (TVD) between the hardware-sampled probability distribution and
the ideal statevector distribution, for both raw and ZNE-corrected results.

% --- AUTO-GENERATED TABLE: paste output of cell 6 in hardware_ibm_colab.ipynb ---
\begin{table}[h]
  \centering
  \caption{QOS oracle circuit fidelity on IBM [BACKEND\_NAME].}
  \label{tab:hardware}
  \begin{tabular}{cccc}
    \toprule
    $n$ & Fidelity (ZNE) & TVD (ZNE) & TVD (raw) \\
    \midrule
    % PASTE LATEX TABLE FROM NOTEBOOK CELL 8 HERE
    \bottomrule
  \end{tabular}
\end{table}

\subsection*{B.4 Circuit Depth Scaling}

Figure~\ref{fig:circuit_depth} shows the number of two-qubit gates
required to implement the QOS oracle as a function of $n$, for both
the worst-case dense decomposition and the $K$-sparse decomposition
($K = 10\%$ of $N = 2^n$).
The dense decomposition scales as $O(2^n)$ two-qubit gates, while the
sparse decomposition scales empirically as $O(K \log K)$, consistent
with Contribution~1 (adaptive sparse oracle).

% --- PASTE FIGURE FROM circuit_depth_scaling.ipynb ---
\begin{figure}[h]
  \centering
  \includegraphics[width=0.9\textwidth]{figs/circuit_depth_vs_n.pdf}
  \caption{Two-qubit gate count vs circuit size $n$ for dense (all-$K$) and
    sparse ($K = 10\%N$) QOS oracle decompositions.
    The dense circuit scales exponentially; the sparse circuit scales
    sub-linearly in $N$.}
  \label{fig:circuit_depth}
\end{figure}

\subsection*{B.5 Discussion}

At $n = 5$ qubits, we observe a TVD of [TVD\_ZNE\_N5] after ZNE correction
(vs.\ [TVD\_RAW\_N5] without correction), representing a
[ZNE\_IMPROVEMENT]\% improvement in distributional accuracy.
This is consistent with the expected NISQ noise floor for
[GATE\_COUNT\_N5] two-qubit gates on superconducting hardware with
average gate fidelity $\approx [GATE\_FIDELITY]$\%~\cite{bharti2022noisy}.

The hardware results confirm that QOS oracle circuits are compatible
with near-term quantum hardware at the circuit sizes relevant for the
experiments in Section~6, and motivate the use of error mitigation
for larger-scale deployments.

% ============================================================
% References to add to main bibliography:
% \bibitem{ibmquantum2024} IBM Quantum, \url{https://quantum.ibm.com}, 2024.
% \bibitem{temme2017error} Temme et al., Phys. Rev. Lett. 119, 180509 (2017).
% \bibitem{li2017efficient} Li & Benjamin, Phys. Rev. X 7, 021050 (2017).
% \bibitem{bharti2022noisy} Bharti et al., Rev. Mod. Phys. 94, 015004 (2022).
% ============================================================
